\begin{abstract}
Alternating least squares (ALS) is a standard method for training recommendation models, yet GPU implementations run it in FP32 on CUDA cores and leave tensor cores idle, because rating matrices are sparse while tensor cores require dense low-precision tiles. We present APR-BALS, an adaptive-precision blocked ALS system that routes users and items by rating count. Dense entities, which receive more than 99.7\% of the training ratings, assemble their Gram matrices on tensor cores with FP16 operands and FP32 accumulation; the sparse tail remains on a scalar FP32 path, and a blocked batched Cholesky solver keeps the solve stage off the critical path. Against an FP32 baseline sharing its binary, solver, and convergence rule, APR-BALS trains Netflix Prize data 4.1--14.8\texttimes{} faster on an RTX 3060 as the latent dimension grows from 16 to 96, and up to 11.6\texttimes{} faster on MovieLens, while held-out RMSE changes by less than 0.005\%. A Tesla T4 reproduces the trend. A three-way ablation shows that reduced precision contributes at most 1.32\texttimes{}; the gain comes from the tensor-core kernels. APR-BALS also runs 3.0--6.4\texttimes{} faster per iteration than cuMF\_als on identical data splits. All measurements are on NVIDIA Ampere and Turing GPUs.
\end{abstract}
